\documentclass[11pt,a4paper]{article}

\usepackage[margin=2.2cm]{geometry}
\usepackage{microtype}
\usepackage{parskip}
\usepackage[hidelinks]{hyperref}

\setlength{\parindent}{0pt}

\begin{document}

\begin{center}
    {\LARGE\bfseries Executive Summary\par}
    \vspace{0.35cm}
    {\large Tracing Sparse Causal Features in an Open Language Model\par}
    \vspace{0.25cm}
    {\normalsize Evelyn Paskhaver \hfill Word Count: 740\hfill July 12, 2026\par}
\end{center}

\section*{Purpose and importance}

Large language models can produce correct answers without exposing the internal computations that generate them. Mechanistic interpretability identifies internal variables and computational pathways that causally support model behaviour. Anthropic's \emph{On the Biology of a Large Language Model} combines sparse feature decompositions, attribution graphs and interventions to analyse such pathways. Independent reproduction tests whether the resulting graph-derived hypotheses remain valid under causal intervention.

This study implements that workflow in Qwen3-4B-Instruct-2507, an open-weight model with four billion parameters. It examines decimal addition, with emphasis on carry propagation; multi-step reasoning through capital-city questions; and a physics-based reasoning extension using SI-unit questions. The central test is whether compact learned feature panels produce predicted, concept-specific effects under matched controls and held-out confirmation.

\section*{Approach}

The pipeline generated synthetic prompts, recorded multilayer-perceptron activations at seven transformer layers, trained one sparse autoencoder (SAE) per layer, constructed contrastive attribution graphs, and performed feature inhibition and activation swap-in experiments. Complete SAE codes and raw model activations were patched as positive controls. Clean-baseline checks restricted evaluation to prompts exhibiting the expected behaviour.

Reconstruction and sparsity diagnostics characterised the initial weak feature interventions. The first SAEs explained 90.7-95.9\% of held-out activation variance while activating 2,167-3,545 of 8,192 latents per example. High reconstruction fidelity therefore coexisted with a dense representation. A fixed TopK comparison subsequently tested whether direct control of active support improved causal specificity. Candidate models were selected from reconstruction diagnostics before intervention outcomes were examined.

\newpage
\section*{Key findings}

Broad activation patches established that the selected model components contained causal answer information. In arithmetic, replacing target MLP states with matched source states transferred the source answer digit. A source with a different answer digit transferred that digit as well, identifying the broad effect as answer-state transfer rather than carry-specific control. TopK retraining increased sparse intervention magnitude, while graph-derived and predeclared ten-feature panels remained non-specific relative to matched no-carry controls. An output-digit-conditioned analysis then ranked all 57,344 SAE latents on separate discovery cases. The ten-feature primary was null; a secondary 20-feature panel was frozen and tested once on 32 intervention-untouched cases. Inhibition changed the carry-target gap by $-0.059$ logits and the no-carry control by $+0.027$, producing a paired effect of $-0.086$ (bootstrap 95\% interval $-0.137$ to $-0.039$). This establishes a replicated, distributed carry-associated effect at the logit level. Capital-city swaps produced limited transfer even after high-confidence prompts replaced the original lower-confidence examples.

The strongest effect occurred for SI units. The selected TopK SAE retained 93.6\% mean held-out variance with approximately 100 active latents, 35.5 times fewer than the original units SAE. Features from one force attribution graph were ranked on eight discovery systems, and the resulting ten-feature panel was frozen before evaluation on 16 disjoint physical systems. Swapping force-source values into energy prompts increased the newtons-versus-joules logit gap by 1.234 logits (bootstrap 95\% interval 1.148--1.332). A matched mass-source control changed the gap by $-0.090$ ($-0.129$ to $-0.051$), giving a paired specificity effect of 1.324 logits (1.223--1.430). Every confirmation case moved more strongly under the force source than under the mass control. Eight of the ten latents occurred at layer 28, and the panel reproduced 97.5\% of the force-source shift obtained from all 73 positive graph features.

Joules remained the top prediction in every confirmation case. The demonstrated endpoint is therefore selective control of unit-token preference by a compact force-associated panel. The experiment does not test unique necessity or identify a complete circuit.

\section*{Implications and next steps}

The study completes a small-scale reproduction of the sparse-feature, attribution-graph and intervention workflow, with domain-dependent causal results. The force panel met a predeclared criterion. The smaller arithmetic effect emerged from a secondary analysis and was subsequently confirmed through a frozen, one-shot replication. Capital-city transfer remained limited. Together, these outcomes show that reconstruction fidelity, feature decodability and causal specificity are distinct properties. The positive interpretations follow from matched controls, frozen feature sets and held-out evaluation.

The next experimental priority is replication of both selection procedures across independently trained SAE seeds, using larger frozen confirmation pools and more random feature panels. Arithmetic SAEs should be trained at the teacher-forced decimal positions used for intervention, with carry status and answer digit varied independently. Per-layer or a small cross-layer transcoder with explicit error nodes would provide a closer architectural comparison with Anthropic's method.

All principal code, configurations, notebooks and output formats are documented in the public repository at
\url{https://gitlab.developers.cam.ac.uk/phy/data-intensive-science-mphil/assessments/projects/ep787}. The reproducibility pack allows the positive units result, the controlled negative results, and the methodological diagnostics to be audited independently.

\end{document} 
